\subsection{标量本质可积函数}

若 \(\mathbf{f}\) 是标量本质 \(\mu\) -可积映射（由 \(\mathrm{T}\) 到 \(\mathrm{F}\) 中），则映射 \(\mathbf{z}' \mapsto \int \langle \mathbf{f}(t), \mathbf{z}' \rangle d\mu(t)\) 是 \(\mathrm{F}'\) 上的线性形式，即代数对偶 \(\mathrm{F}'^*\) 的一个元素。

定义 1. --- 我们把 f 关于 \(\mu\) 的积分记作 \(\int f d\mu\) 或 \(\int f(t) d\mu(t)\)，把它定义为下式给出的 \(F'^{*}\) 中的元素：

\[
\left\langle \mathbf {z} ^ {\prime}, \int \mathbf {f} d \mu \right\rangle = \int \left\langle \mathbf {z} ^ {\prime}, \mathbf {f} \right\rangle d \mu
\]

对一切 \(\mathbf{z}' \in \mathrm{F}'\)。

若 f 连续且具有紧支集，则它标量可积，且定义 1 与 Ch. III, §3, No.~1 中给出的 f 的积分定义一致。另一方面，若 F 是 Banach 空间且 f 本质可积（Ch. V, §1, No.~3, 定义 3），则 f 标量本质可积，且定义 1 与第 V 章中给出的 f 的积分定义一致（Ch. V, §1, No.~3 与 Ch. IV, §4, No.~2, 定理 1 的推论 1）。

例. --- 设 X 为局部紧空间，\(t \mapsto \lambda_{t}\) 为 T 到 X 上测度空间 \(\mathcal{M}(X)\) 中的映射。说族 \(t \mapsto \lambda_{t}\) 是 \(\mu\) -适宜的，就是指它由正测度组成，且映射 \(t \mapsto \lambda_{t}\) 是标量本质 \(\mu\) -可积的与 \(\mu\) -可测的（关于拓扑 \(\sigma(\mathcal{M}(X), \mathcal{K}(X))\)）。² 它关于 \(\mu\) 的积分就是在 Ch. V, §3, No.~1 中记作 \(\int \lambda_t d\mu(t)\) 的那个测度。

注. --- 1) 若 F 是有限维的，则 T 到 F 中的每个标量本质可积映射都是本质可积的（Ch. V, §1, No.~3）。然而在一般情形，紧空间 T 上的标量可忽略函数甚至可能不是 \(\mu\) -可测的（习题 12）。

\begin{enumerate}
\def\labelenumi{\arabic{enumi})}
\setcounter{enumi}{1}
\tightlist
\item
  显然，\(\mathbf{f}\) 的积分只依赖于 \(\mathbf{f}\) 模 \(\mathbf{T}\) 到 \(\mathbf{F}\) 中的标量局部 \(\mu\) -可忽略映射所成空间的类。注意，标量局部可忽略函数 \(\mathbf{g}\) 不一定几乎处处为零（习题 12）。然而，当 \(\mathbf{F}'\) 中存在一个序列 \((\mathbf{z}_n')\)（关于拓扑 \(\sigma(\mathbf{F}',\mathbf{F})\) 稠密）时，情形确实如此：因为若 \(\mathbf{H}_n\) 是 \(t \in \mathbf{T}\) 中使 \(\langle \mathbf{g}(t), \mathbf{z}_n' \rangle \neq 0\) 的点所成的局部可忽略集，则 \(\mathbf{H}\)（诸 \(\mathbf{H}_n\) 之并）局部可忽略，且对每个 \(t \notin \mathbf{H}\)，有 \(\langle \mathbf{g}(t), \mathbf{z}_n' \rangle = 0\)（对一切 \(n\)），从而 \(\mathbf{g}(t) = 0\)。
\end{enumerate}

设 \(u\) 是 \(F\) 到 Hausdorff 局部凸空间 \(G\) 中的连续线性映射；它的转置 \(^{t}u\) 是 \(G'\) 到 \(F'\) 中的线性映射，而（代数）转置 \(^{t}(^{t}u)\) 是 \(F'^{*}\) 到 \(G'^{*}\) 中的线性映射，它是 \(u\) 的扩张，我们仍把它记作 \(u\)。按此约定：

命题 1. --- 若 f 是 T 到 F 中的标量本质 μ-可积映射，则映射 u ∘ f 是标量本质 μ-可积的，且

\[
\int (u \circ \mathbf {f}) d \mu = u \Big (\int \mathbf {f} d \mu \Big).
\]

对每个 \(z^{\prime} \in G^{\prime}\)，\(\langle z^{\prime}, u \circ f \rangle = \langle {}^{t} u(z^{\prime}), f \rangle\)，由此得第一个结论；第二个结论由公式

\[
\left\langle \mathbf {z} ^ {\prime}, \int (u \circ \mathbf {f}) d \mu \right\rangle = \int \left\langle \mathbf {z} ^ {\prime}, u \circ \mathbf {f} \right\rangle d \mu = \left\langle {} ^ {t} u (\mathbf {z} ^ {\prime}), \int \mathbf {f} d \mu \right\rangle = \left\langle \mathbf {z} ^ {\prime}, u \left(\int \mathbf {f} d \mu\right) \right\rangle .
\]

特别地，若 f 是标量本质 \(\mu\) -可积的，则当 F 的拓扑换成更粗的拓扑时，它仍是标量本质 \(\mu\) -可积的。

命题 2. --- 设 f 是 T 到 F 中的标量本质 \(\mu\) -可积映射。对每个数值函数 \(g \geqslant 0\)（\(\mu\) -可测且有界），T 到 F 中的映射 \(t \mapsto g(t)\mathbf{f}(t)\)（记作 gf 或 fg）是标量本质 \(\mu\) -可积的，f 是标量本质 \((g \cdot \mu)\) -可积的，且

\[
\int \mathbf {f} d (g \cdot \mu) = \int \mathbf {f} g d \mu .
\]

这是公式 \(\langle z', gf\rangle = g\langle z', f\rangle\)（对一切 \(z' \in F'\)）以及公式 \(\int h d(g \cdot \mu) = \int hg d\mu\)（对每个本质 \(\mu\) -可积的数值函数 h）（Ch. V, §5, No.~3, 定理 1）的直接推论。

关于本质可积数值函数的大量命题可以逐字逐句地改述为关于标量本质可积向量值函数的命题。其中较重要者，我们提请注意函数关于由密度定义的测度（Ch. V, §5, No.~3, 定理 1）、关于测度的像（Ch. V, §6, No.~2, 定理 1）、关于诱导测度（Ch. V, §7, No.~1, 定理 1）、或关于正测度可和族之和（Ch. V, §2, No.~2, 命题 1 与 3 及命题 1 的推论 3）本质可积的条件。这些改述留给读者。

然而，为了得到与“二重积分”定理（Ch. V, §3, No.~3, 定理 1 与 §8, No.~4, 定理 1（Lebesgue--Fubini 定理））相应的命题，必须加强假设（参见习题 1）；把前面引用的定理应用于每个函数 \(\langle z', f \rangle\)（其中 \(z' \in F'\)），就得到下列命题：

命题 3. --- 设 X 为局部紧空间，\(t \mapsto \lambda_{t}\) 为 X 上正测度的 \(\mu\) -适宜族 \(^{3}\)（Ch. V, §3, No.~1, 定义 1），且 \(\nu = \int \lambda_{t} d\mu(t)\)。设 f 为 X 到 F 中的映射；假设：\(1^{\circ} f\) 是标量 \(\nu\) -可积的；\(2^{\circ}\) 存在局部 \(\mu\) -可忽略集 \(N \subset T\)，使对每个 \(t \notin N\)，f 是标量 \(\lambda_{t}\) -可积的且 \(\int f d\lambda_{t} \in F\)。那么，函数 \(t \mapsto \int f d\lambda_{t}\)（对 \(t \notin N\) 定义）是标量本质 \(\mu\) -可积的，且 \(^{4}\)

\[
\int \mathbf {f} (x) d \nu (x) = \int d \mu (t) \int \mathbf {f} (x) d \lambda_ {t} (x).
\]

命题 4. --- 设 T 与 T′ 为两个局部紧空间，\(\mu\)（相应地 \(\mu'\)）为 T（相应地 T′）上的正测度，\(\nu = \mu \otimes \mu'\) 为 X = T × T′ 上的乘积测度。设 f 为 X 到 F 中的映射。假设：\(1^{\circ} f\) 是标量 \(\nu\) -可积的；\(2^{\circ}\) 存在局部 \(\mu\) -可忽略集 N ⊂ T，使对每个 \(t \notin N\)，映射 \(t' \mapsto f(t, t')\) 是标量 \(\mu'\) -可积的，且 \(\int f(t, t') d\mu'(t') \in F\)。那么，函数 \(t \mapsto \int f(t, t') d\mu'(t')\)（对 \(t \notin N\) 定义）是标量本质 \(\mu\) -可积的，且 \(^{4}\)

\[
\iint \mathbf {f} (t, t ^ {\prime}) d \mu (t) d \mu^ {\prime} (t ^ {\prime}) = \int d \mu (t) \int \mathbf {f} (t, t ^ {\prime}) d \mu^ {\prime} (t ^ {\prime}).
\]

